% Generated by scripts/assemble.py; corrections belong in source/corrections.json.
\examhead{168}{电子科技大学 2025 年数学分析试卷}{159}
% p160-b27; source PDF page 160
\par\Needspace{5\baselineskip}
\blockanchor{p160-b27}
\majorquestion{一}填空题. 每题 5 分, 共 30 分.\par
% p160-b28; source PDF page 160
\blockanchor{p160-b28}
\question{1}求 $\lim_{n\to \infty}n\left[\left(1 + \frac{1}{n}\right)^n -e\right] =$\par
% p160-b29; source PDF page 160
\blockanchor{p160-b29}
\question{2}已知 $z = x \ln \frac{x}{y}$, 求 $\mathrm{d}^2 z \Big|_{(1,1)} =$ \_\_\_\_.\par
% p161-b00; source PDF page 161
\blockanchor{p161-b00}
\question{3}计算 $\oint_{L} (xy + yz + zx) \, \mathrm{d}s =$ \_\_\_\_，其中 $L$ 为 $x^{2} + y^{2} + z^{2} = a^{2}$ 与 $x + y + z = 0$ 的交线.\par
% p161-b01; source PDF page 161
\blockanchor{p161-b01}
\question{4}幂级数  $\sum_{n=1}^{\infty}\left(1+\frac{1}{2}+\cdots+\frac{1}{n}\right)^{2}x^{n}$  的收敛域为 \_\_\_\_.\par
% p161-b02; source PDF page 161
\blockanchor{p161-b02}
\question{5}已知 $f(x) = x(\pi - x), x \in [0, \pi]$，展开为余弦级数为 \_\_\_\_.\par
% p161-b03; source PDF page 161
\blockanchor{p161-b03}
\question{6}已知 $b > a > 0$, 计算 $\int_0^1 \frac{x^b - x^a}{\ln x} \, \mathrm{d}x =$ \_\_\_\_.\par
% p161-b04; source PDF page 161
\par\Needspace{5\baselineskip}
\blockanchor{p161-b04}
\majorquestion{二}计算题. 每题 10 分, 共 50 分.\par
% p161-b05; source PDF page 161
\blockanchor{p161-b05}
\question{7}已知数列 $\{a_{n}\}: a_{1} = 0, a_{2m} = \frac{a_{2m - 1}}{2}, a_{2m + 1} = \frac{1}{2} + a_{2m}$, 求 $\{a_{n}\}$ 的上下极限.\par
% p161-b06; source PDF page 161
\blockanchor{p161-b06}
\question{8}求 $\int_0^{\frac{\pi}{2}}\frac{f(x)}{\sqrt{x}}\mathrm{d}x,$ 其中 $f(x) = \int_{\sqrt{\frac{\pi}{2}}}^{\sqrt{x}}\frac{1}{1 + \tan u^2}\mathrm{d}u.$\par
% p161-b07; source PDF page 161
\blockanchor{p161-b07}
\question{9}计算 $\lim_{n\to \infty}\frac{1}{n}\sum_{j = 1}^{n}\sum_{k = 1}^{n}\frac{j}{j^2 + k^2}.$\par
% p161-b08; source PDF page 161
\blockanchor{p161-b08}
\question{10}计算 $\int_{L} \frac{y \, \mathrm{d}x - (x - 1) \, \mathrm{d}y}{(x - 1)^2 + y^2}$, 其中 $L$ 为椭圆 $x^2 - 4x + 4y^2 = 0$, 取逆时针方向.\par
% p161-b09; source PDF page 161
\blockanchor{p161-b09}
\question{11}求第二类曲面积分 $\iint_{S} \frac{x}{x^2 + y^2 + z^2} \, \mathrm{d}y \, \mathrm{d}z$, 其中 $S$ 为圆柱面 $x^2 + y^2 = 1 (-1 \leq z \leq 1)$, 方向为外侧.\par
% p161-b10; source PDF page 161
\par\Needspace{5\baselineskip}
\blockanchor{p161-b10}
\majorquestion{三}证明题. 每题10分, 共40分.\par
% p161-b11; source PDF page 161
\blockanchor{p161-b11}
\question{12}$A$ 是由数码0,1组成的所有数列的集合, 证明 $A$ 不可数.\par
% p161-b12; source PDF page 161
\blockanchor{p161-b12}
\question{13}证明函数 $f(x) = \sqrt{x} \ln x$ 在 $[1, +\infty)$ 上一致连续.\par
% p161-b13; source PDF page 161
\blockanchor{p161-b13}
\question{14}证明级数 $\sum_{n=1}^{\infty} \frac{\sin(n+1) \cos(n-1)}{n^p}$ 在 $p > 1$ 时收敛, 在 $p \leq 1$ 时发散.\par
% p161-b14; source PDF page 161
\blockanchor{p161-b14}
\question{15}证明级数 $\sum_{n=1}^{\infty} (-1)^{n} \frac{x^{2} + n}{n^{2}}$ 在闭区间 $[a, b]$ 上一致收敛, 但对任意的 $x$ 不绝对一致收敛.\par
% p161-b15; source PDF page 161
\par\Needspace{5\baselineskip}
\blockanchor{p161-b15}
\majorquestion{四}综合题. 每题 15 分, 共 30 分.\par
% p161-b16; source PDF page 161
\blockanchor{p161-b16}
\question{16}(x) 表示 $x$ 的小数部分, 已知 $f(x) = \sum_{n=1}^{\infty} \frac{(nx)}{n^2} (x \in \mathbb{R})$.\par
% p161-b17; source PDF page 161
\blockanchor{p161-b17}
(1) 求 $f(x)$ 所有间断点.\par
% p161-b18; source PDF page 161
\blockanchor{p161-b18}
(2) 证明 $f(x)$ 在有界闭区间 $[a, b]$ 上黎曼可积.\par
% p161-b19; source PDF page 161
\blockanchor{p161-b19}
\question{17}已知数列 $\{a_{n}\}$ 满足 $a_{n+1} > a_{n} > 0$, 且 $\sum_{n=1}^{\infty} a_{n}$ 收敛.\par
% p161-b20; source PDF page 161
\blockanchor{p161-b20}
(1) 证明: $\lim_{n\to \infty}na_n = 0.$\par
% p161-b21; source PDF page 161
\blockanchor{p161-b21}
(2) 证明: $\sum_{n=1}^{\infty} n(a_n - a_{n+1})$ 收敛.\par
